\documentclass[11pt]{article}
\usepackage[T1]{fontenc}
\usepackage[utf8]{inputenc}
\usepackage{lmodern}
\usepackage[letterpaper,margin=0.82in]{geometry}
\usepackage{microtype}
\usepackage{setspace}
\usepackage{booktabs}
\usepackage{longtable}
\usepackage{tabularx}
\usepackage{array}
\usepackage{enumitem}
\usepackage{xcolor}
\usepackage{amsmath}
\usepackage{fancyhdr}
\usepackage{titlesec}
\usepackage{url}
\usepackage[hidelinks,bookmarksopen=true]{hyperref}

\definecolor{kitblue}{RGB}{28,74,110}
\definecolor{kitgray}{RGB}{90,96,102}
\definecolor{lightgray}{RGB}{245,246,247}
\definecolor{green}{RGB}{27,105,65}
\definecolor{orange}{RGB}{148,91,0}
\definecolor{red}{RGB}{140,40,40}

\hypersetup{
  pdftitle={EvalAudit Paper Strategy and Execution Plan},
  pdfauthor={Edward Wang},
  pdfsubject={HELM reproducibility, EvalAudit, technical paper strategy, and small-model evaluation},
  pdfkeywords={EvalAudit, HELM, reproducibility, execution provenance, Qwen, small models}
}

\setstretch{1.04}
\setlength{\parindent}{0pt}
\setlength{\parskip}{5pt}
\setlist[itemize]{leftmargin=1.4em,itemsep=2pt,topsep=3pt}
\setlist[enumerate]{leftmargin=1.6em,itemsep=2pt,topsep=3pt}
\renewcommand{\arraystretch}{1.15}
\newcolumntype{Y}{>{\raggedright\arraybackslash}X}
\newcolumntype{P}[1]{>{\raggedright\arraybackslash}p{#1}}
\newcommand{\code}[1]{\texttt{\detokenize{#1}}}
\newcommand{\ready}{\textcolor{green}{\textbf{Ready}}}
\newcommand{\partstat}{\textcolor{orange}{\textbf{Partial}}}
\newcommand{\blocked}{\textcolor{red}{\textbf{Blocked}}}
\newcommand{\callout}[2]{%
  \begin{center}\fcolorbox{kitgray}{lightgray}{\begin{minipage}{0.93\textwidth}
  \textbf{#1}\quad #2
  \end{minipage}}\end{center}}

\pagestyle{fancy}
\fancyhf{}
\fancyhead[L]{\small EvalAudit paper strategy}
\fancyhead[R]{\small Revised July 15, 2026}
\fancyfoot[C]{\thepage}
\titleformat{\section}{\Large\bfseries\color{kitblue}}{\thesection}{0.7em}{}
\titleformat{\subsection}{\large\bfseries\color{kitgray}}{\thesubsection}{0.7em}{}

\begin{document}
\pagenumbering{gobble}

\begin{titlepage}
\centering
\vspace*{0.85in}
{\Huge\bfseries\color{kitblue} EvalAudit Paper Strategy\par}
\vspace{0.18in}
{\LARGE Framing, Scope, Evidence, and a Three-Week Execution Plan\par}
\vspace{0.55in}
{\Large Edward Wang --- Kitware Internship 2026\par}
\vfill
\begin{minipage}{0.84\textwidth}
\small
This reference consolidates the internship chronology, status slides, repository history, journals, runbooks, the inherited Every Eval Ever paper, the TMLR draft skeleton, the OLMo deployment-match confirmation path, and the subsequent paper-framing discussion. It distinguishes what can be published now from what requires additional empirical work.

\textbf{Evidence cutoff:} repository snapshot through July 14, 2026; interpretive revision July 15, 2026.
\end{minipage}
\end{titlepage}

\pagenumbering{roman}
\tableofcontents
\clearpage
\pagenumbering{arabic}

\section{Executive recommendation}

With approximately 2.5--3 weeks remaining, the best target is neither a pure position paper nor a broad benchmark paper. It is:

\begin{quote}
\textbf{A systems-and-measurement paper with a clear position, one deep modern empirical case, one historical forensic case, and an optional narrow current-model demonstration.}
\end{quote}

The paper must remain publishable if the remaining GPU work does not complete. Its core contribution is EvalAudit and a layered account of benchmark reproduction. Its empirical upgrade is an ordinary-path OLMo confirmation. A Qwen comparison is valuable only after that core result and artifact are secured. Fable-distilled models should be treated as a follow-on unless their data and training provenance are already complete.

\callout{Primary decision}{Secure a correct, auditable manuscript first. Increase model count only after the ordinary-HELM OLMo confirmation, result regeneration, and evidence packaging are complete.}

\section{The two-paper lineage}

The repository contains two scientifically distinct paper lines.

\subsection{Inherited EEE paper}

The NeurIPS Every Eval Ever (EEE) work uses a normalized representation to detect discrepancies on Pythia-6.9B, Vicuna-7B-v1.3, and Falcon-7B across fourteen HELM benchmarks. Its scientific endpoint is discrepancy detection and forensic inspection. Its appendix explicitly leaves unresolved whether residual differences arise from checkpoint, quantization, precision, tokenizer, or generation-kernel choices.

This is inherited, substantially pre-internship work. It should be cited and used as motivation, not retold as a new contribution.

\subsection{Internship follow-on}

The pre-existing TMLR skeleton is the natural home for the internship's distinct question:

\begin{quote}
How far can execution-provenance reconstruction attribute or close the discrepancies that normalized comparison reveals, and when is the original measurement no longer identifiable?
\end{quote}

The clean positioning sentence is:

\callout{Paper relationship}{EEE detects reproducibility gaps; the follow-on paper studies their attribution, possible closure, and identifiability.}

\section{Recommended thesis and paper identity}

\subsection{Defensible core thesis}

\begin{quote}
\textbf{Public LLM benchmark artifacts often identify the evaluation recipe but not the complete execution instrument. Exact run-spec replay and controlled reconstruction of model, tokenizer, precision, engine, and harness provenance can attribute---and sometimes close---apparent reproducibility gaps. When that provenance has been lost or post-hoc transformed, the original measurement may remain non-identifiable.}
\end{quote}

This thesis is supported even if the ordinary-path OLMo confirmation fails.

\subsection{Forward-looking systems motivation}

The strongest ``why care'' framing is not old-model preservation for its own sake. Historical HELM runs provide forensic evidence from which to derive requirements for trustworthy current evaluation.

\begin{quote}
\textbf{Auditable evaluation is necessary to determine whether measured improvements in small open models exceed variation introduced by the evaluation substrate itself.}
\end{quote}

The small-model story should be the application of EvalAudit, not a separate benchmark paper attached to a HELM postmortem.

\subsection{Recommended title}

\begin{quote}
\textbf{EvalAudit: Reconstructing the Execution Substrate of LLM Benchmarks}
\end{quote}

Possible subtitle:

\begin{quote}
From HELM reproduction failures to auditable evaluation of current local models
\end{quote}

Alternative conceptual title:

\begin{quote}
\textbf{A Benchmark Recipe Is Not an Experiment: Attribution and Identifiability in LLM Evaluation}
\end{quote}

\section{Current evidence status}

\begin{longtable}{@{}P{0.29\textwidth}P{0.13\textwidth}P{0.49\textwidth}@{}}
\toprule
Evidence / capability & Status & Defensible interpretation \\
\midrule
Exact public \code{run_spec.json} replay & \ready & Implemented and shown to correct visible recipe drift, including the documented BBQ prompt mismatch. \\
Containerized HELM execution and GPU leasing & \ready & Reusable infrastructure exists; individual family grids still require result-level verification. \\
OLMo tokenizer and template diagnoses & \ready & Special-token and chat-template behavior can change the actual model stimulus and outputs. \\
OLMo float32/engine result & \partstat & Highest-agreement candidates in an oracle-scored 12-instance IFEval probe; not an ordinary HELM reproduction or proof of the historical deployment. \\
OLMo Hugging Face dense-model divergence & \partstat & Prompt mismatch was rejected for an inspected case; root cause remains unresolved and scientifically relevant. \\
Classic HELM source archaeology & \ready & Strong evidence of carried-forward artifacts, migrated class paths, and an inferred compatible unreleased source interval. \\
Classic proxy execution & \partstat & Era machinery exists; full original-model proxy comparisons were not established in the committed runbook state. \\
GPT-OSS workflow & \partstat & Null-content compatibility issue diagnosed; exact-path result freshness and acceptance status require reconciliation. \\
Qwen workflow & \partstat & Combined fan-out is implemented; weight mappings and complete GPU acceptance require verification. \\
Population-level reproducibility rate & \blocked & No frozen denominator and no fully regenerated publication store. \\
Broad small-model progress claim & \blocked & Requires exact model revisions, fair protocols, substrate sensitivity, and complete benchmark results. \\
\bottomrule
\end{longtable}

\section{The technical model}

Represent a measured benchmark result as

\[
Y = F(R,D,I,A,U),
\]

where:

\begin{itemize}
  \item $R$: recorded recipe---scenario, data, prompt construction, decoding, scorer;
  \item $D$: deployment---weights, model/tokenizer revisions, template, dtype, quantization;
  \item $I$: execution instrument---harness, engine, packages, kernels, hardware, topology, batching;
  \item $A$: artifact interpretation---schema, migration, canonicalization, carried-forward outputs;
  \item $U$: residual stochastic and numerical influences.
\end{itemize}

The public record exposes only a projection of these variables. Reproduction is identifiable at tolerance $\epsilon$ when all materially compatible configurations produce equivalent results under a declared output-, metric-, or conclusion-level criterion.

Recommended outcome classes:

\begin{enumerate}
  \item Recipe mismatch
  \item Deployment mismatch
  \item Execution-instrument mismatch
  \item Artifact-interpretation mismatch
  \item Residual disagreement under controlled equivalence
  \item Historically non-identifiable reproduction
  \item Output-level reproduction
  \item Metric-level reproduction
  \item Robust conclusion despite non-identical outputs
\end{enumerate}

\section{EvalAudit as the systems contribution}

EvalAudit should be described as an evaluation auditing and experimental-control system, not merely a HELM wrapper.

\subsection{Evaluation capsules}

Each run should preserve:

\begin{itemize}
  \item frozen benchmark specification, harness commit, and container digest;
  \item dataset and scorer versions;
  \item model and tokenizer repository revisions;
  \item chat-template text or hash and tokenized-prompt hash;
  \item dtype, quantization, engine, attention implementation, and generation configuration;
  \item batching, cache, tensor-parallel, numerical, GPU, CUDA, and driver settings;
  \item raw requests, token IDs, responses, metrics, and artifact-transformation history.
\end{itemize}

\subsection{Layered differential comparison}

EvalAudit should identify the earliest divergence layer:

\[
\text{recipe} \rightarrow \text{prompt} \rightarrow \text{tokens} \rightarrow
\text{generation} \rightarrow \text{text} \rightarrow \text{metric}.
\]

\subsection{Controlled replay}

The system should support counterfactual variation of one component while freezing the others: revision, tokenizer, template, special tokens, dtype, engine, harness era, and---where feasible---hardware-sensitive settings.

\subsection{Structured reports}

Reports should separate coverage and failure classes rather than emitting one undifferentiated reproduction rate.

\section{The blocking empirical result}

The highest-priority experiment is to move the OLMo candidate from the deployment probe into the real HELM path.

\begin{enumerate}
  \item Freeze discovery instances and candidate configurations.
  \item Determine whether the winning candidate requires probe-only \code{add_special_tokens=False} behavior.
  \item Implement required behavior through the actual tokenizer/client path.
  \item Execute the exact public \code{run_spec.json} through normal HELM.
  \item Generate ordinary HELM artifacts.
  \item Compare through the normal pair-report pipeline.
  \item Evaluate on held-out instances and preferably another benchmark.
  \item Preserve the full candidate surface, not only the winner.
\end{enumerate}

A success supports a concrete attribution-and-closure claim. A failure remains publishable: it demonstrates that output-matching probes do not automatically translate into faithful benchmark execution.

\section{Minimum empirical package}

\begin{longtable}{@{}P{0.24\textwidth}P{0.14\textwidth}P{0.14\textwidth}P{0.39\textwidth}@{}}
\toprule
Study & Models & Benchmarks & Purpose \\
\midrule
OLMo ordinary-path confirmation & 2 & 2 & Dense plus MoE; short deterministic plus IFEval; naive, exact recipe, reconstructed deployment. \\
OLMo bounded ablation & 1--2 & 1--2 & Dtype, template/generation prompt, and special-token behavior. \\
Repeatability baseline & 1 & 1 & Three same-machine runs; cross-machine only if operational quickly. \\
Historical case & 1 & 1--2 & Artifact/source proof; optionally two proxy-era executions. \\
Qwen demonstration & 2 & 2--3 & Compare model-generation delta with variation under one alternative plausible substrate. \\
\bottomrule
\end{longtable}

Outcome metrics should include exact tokens, normalized text, first divergence position, common-prefix length, task-specific correctness or IFEval compliance, aggregate score difference, and ranking stability.

\section{Small-model application}

\subsection{Scientific question}

The future-facing question is:

\begin{quote}
Is the new or fine-tuned model genuinely better, or is the measured difference an artifact of the evaluation configuration?
\end{quote}

For a model pair, compare

\[
\Delta_{\mathrm{model}} = Y_{\mathrm{new}} - Y_{\mathrm{old}}
\]

with

\[
\Delta_{\mathrm{substrate}} = \max_c Y_{m,c} - \min_c Y_{m,c}.
\]

A robust progress claim should exceed or remain stable across the tested substrate variation.

\subsection{Qwen scope}

Within the internship, use at most one carefully matched pair, exact model and tokenizer revisions, two or three benchmarks that discriminate within the small-model cohort, and a pinned EvalAudit environment. Separate native model templates from standardized evaluation semantics, and do not mix thinking/reasoning modes with direct-answer modes without a matched compute or token budget.

The purpose is to demonstrate EvalAudit's relevance, not to establish a comprehensive Qwen leaderboard.

\subsection{Fable-distilled models}

Fable-tuned models become a headline result only if the following already exist:

\begin{itemize}
  \item exact parent and teacher checkpoints;
  \item teacher prompts and sampling settings;
  \item training-data manifest, filters, deduplication, and contamination analysis;
  \item complete training hyperparameters, seeds, and checkpoint-selection rule;
  \item held-out and unrelated-task evaluation partitions.
\end{itemize}

Otherwise, include them as exploratory future work. Adding incomplete synthetic-data experiments would create more reviewer risk than value during the final weeks.

\section{Cost and hardware}

The paper should not describe hardware simply as an unmitigable source of non-reproducibility. Some historical hardware cannot be recreated, but it can be recorded, varied, bounded, and tested for conclusion robustness.

Distinguish:

\begin{enumerate}
  \item \textbf{Exact replay:} same recorded instrument is available.
  \item \textbf{Equivalent replay:} a different instrument produces outputs or metrics within a declared tolerance.
  \item \textbf{Robust conclusion:} exact outputs differ, but the scientific conclusion or ranking is stable across tested instruments.
\end{enumerate}

For current small models, report GPU-hours, peak memory, throughput, latency, and storage/audit overhead. This makes the paper relevant to local-model development and exposes whether a model's gain is worth its resource cost.

\section{Claims and language discipline}

\subsection{Safe claims}

\begin{itemize}
  \item A benchmark recipe does not fully specify the experiment.
  \item Exact-spec replay prevents documented recipe-reconstruction errors.
  \item Unrecorded tokenizer, template, dtype, and serving choices can materially change outputs in the studied cases.
  \item Some missing provenance is experimentally recoverable; some historical provenance is not uniquely recoverable.
  \item EvalAudit provides infrastructure to preserve and diagnose these layers.
\end{itemize}

\subsection{Claims to avoid until further experiments}

\begin{itemize}
  \item Most HELM results are reproducible.
  \item Float32 was definitively the original TogetherAI configuration.
  \item Precision is the dominant hidden variable in general.
  \item Residual gaps are small or almost always attributable.
  \item Qwen 3.x/3.5 definitively advances over Qwen 2.5 across the capability spectrum.
  \item Fable fine-tuning reliably improves small models.
\end{itemize}

Use ``highest-agreement tested candidate'' rather than ``the original deployment,'' and ``consistent with'' rather than ``proved'' when multiple latent configurations remain possible.

\section{Three-week execution plan}

\subsection{Days 1--3: freeze and draft}

\begin{itemize}
  \item Freeze thesis, contribution list, model/benchmark roster, and denominator.
  \item Build a run-status table separating probe, smoke, exact-path, stale, fresh, and completed results.
  \item Regenerate all non-GPU provenance analyses.
  \item Draft introduction, EEE relationship, technical model, EvalAudit system, and historical case.
\end{itemize}

At this point the manuscript should already be viable as a position/experience paper.

\subsection{Days 4--8: OLMo confirmation}

Resolve the request-path blocker and attempt one dense and one MoE model through ordinary HELM on two benchmarks. Regenerate reports with current code.

\callout{Decision gate}{If no ordinary-path confirmation is complete by the end of this block, stop broadening experiments. Preserve the failure as a result and finish the position/systems paper.}

\subsection{Days 9--12: bounded ablation and repeatability}

Run a small staged or $2\times2\times2$ design over dtype, template/generation prompt, and special-token handling. Add three repeats of the final configuration. Prefer held-out confirmation over more factors.

\subsection{Days 13--15: historical evidence and artifact}

Regenerate hashes, source compatibility, and the historical timeline. Run a proxy-era comparison only if the path is already operational. Package immutable evidence and figure scripts.

\subsection{Days 16--18: optional Qwen demonstration}

Proceed only after the paper's core empirical and artifact milestones are complete. Use one matched pair, two or three benchmarks, and one alternative plausible execution configuration.

\subsection{Final days: freeze}

Stop experiments, audit every claim against stored evidence, complete limitations, regenerate every figure from scripts, and archive the paper artifact.

\section{Recommended display items}

The highest-value figures and tables are:

\begin{enumerate}
  \item Layered measurement-instrument diagram.
  \item EEE-to-follow-on paper relationship diagram.
  \item Corpus eligibility and completion funnel.
  \item OLMo agreement waterfall: naive, exact recipe, reconstructed deployment.
  \item OLMo bounded-ablation table or heatmap.
  \item First-divergent-token examples.
  \item Historical HELM source/artifact timeline.
  \item Provenance field table: recorded, missing, material, recoverable, recommended.
  \item Optional Qwen model delta versus substrate variation.
  \item Capability-cost frontier for the small-model demonstration.
\end{enumerate}

\section{Submission-readiness criteria}

The hybrid paper is ready when:

\begin{itemize}
  \item inherited EEE results and internship contributions are explicitly separated;
  \item every result is labeled probe, smoke, exact-path, historical inference, or completed comparison;
  \item all cited tables are regenerated from current code and a frozen denominator;
  \item the ordinary-path OLMo outcome is reported honestly, whether success or failure;
  \item historical inferences are labeled as source-compatibility inferences;
  \item raw evidence, manifests, revisions, environments, and figure scripts are archived;
  \item the paper does not depend on unfinished Qwen or Fable results.
\end{itemize}

A broader empirical paper additionally requires held-out configuration validation, larger benchmark samples, controlled ablations, local-local and cross-machine baselines, completed modern-family breadth, and statistical uncertainty.

\section{Final recommendation}

The highest expected-value strategy is:

\begin{quote}
\textbf{Write the position and system paper now; use one end-to-end OLMo study to raise its technical level; add a narrow Qwen demonstration only if the core work finishes early; reserve the full small-model and Fable study for the expanded follow-on paper.}
\end{quote}

This plan preserves a publishable output under failure, directs scarce GPU time toward the single experiment that changes the paper's evidentiary status, and leaves a coherent path to the higher-value small-model measurement paper.

\end{document}
